\usepackage[utf8]{inputenc}
\usepackage[T1]{fontenc}    % use 8-bit T1 fonts
\usepackage{hyperref}       % hyperlinks
\usepackage{url}            % simple URL typesetting
\usepackage{booktabs}       % professional-quality tables
\usepackage{amsfonts}       % blackboard math symbols
\usepackage{nicefrac}       % compact symbols for 1/2, etc.
\usepackage{microtype}      % microtypography
\usepackage{graphicx}       %image package
\usepackage{bm}
\usepackage{bbm}
%\usepackage{amsmath}
\usepackage{amsthm}
\usepackage{dsfont}
\usepackage{enumitem}
\usepackage{tabularx}
\usepackage[fleqn]{amsmath}
\usepackage[margin=1in]{geometry}
\usepackage{tikz}
\usepackage{listings}
\usepackage{xcolor}
\usepackage{float}
\usepackage[rightcaption]{sidecap}
\usepackage[backend=biber,style=alphabetic,]{biblatex}
\usepackage{titlesec}

\usetikzlibrary{shapes.geometric, arrows}


\DeclareMathOperator{\Tr}{Tr}
\DeclareMathOperator{\Det}{Det}
\DeclareMathOperator{\E}{E}
\DeclareMathOperator{\ER}{ER}
\DeclareMathOperator{\Var}{Var}
\DeclareMathOperator{\argmax}{arg\max}
\DeclareMathOperator{\argmin}{arg\min}
\DeclareMathOperator{\N}{\mathbb{N}}
\DeclareMathOperator{\R}{\mathbb{R}}
\DeclareMathOperator{\ind}{\mathbbm{1}}
\DeclareMathOperator{\I}{I}
\DeclareMathOperator{\Unif}{Unif}
\DeclareMathOperator{\Ber}{Ber}
\DeclareMathOperator{\Binom}{Binom}
\DeclareMathOperator{\Poi}{Poi}
\DeclareMathOperator{\Dir}{Dir}
\DeclareMathOperator{\diag}{diag}
\DeclareMathOperator{\vecop}{Vec}
\DeclareMathOperator{\conv}{Conv}
\DeclareMathOperator{\supp}{Supp}
\DeclareMathOperator{\relu}{ReLU}
\DeclareMathOperator{\asto}{\overset{a.s.}{\longrightarrow}}
\DeclareMathOperator{\pto}{\overset{p}{\longrightarrow}}
\DeclareMathOperator{\dto}{\overset{d}{\longrightarrow}}
\DeclareMathOperator{\hnull}{H_0}
\DeclareMathOperator{\halt}{H_1}


\newcommand{\pderiv}[3][]{\frac{\partial^{#1}}{ \partial #2}\left[ #3 \right]}
\newcommand{\pderif}[3][]{\frac{\partial^{#1} #3}{ \partial #2}}
\newcommand{\s}[2] {\sum_{#1}^{#2}}


\newcommand{\fr}[1]{ \left( #1 \right)}
\newcommand{\fb}[1]{ \left[ #1 \right]}
\newcommand{\fc}[1]{ \left\{ #1 \right\}}
\newcommand{\given}[1]{\left. #1 \right|}
\newcommand{\set}[2]{\left\{#1 \left\vert #2 \right. \right\}}
\newcommand{\norm}[1]{\left\lVert #1 \right\rVert}
\newcommand{\abs}[1]{\left| #1 \right|}

\newcommand{\dls}[1]{{\color{red} #1}}
\newcommand{\bls}[1]{{\color{blue} #1}}


\newcounter{todocounter}
\newcommand{\todonum}[2][]
{\stepcounter{todocounter}\todo[#1]{\textbf{Action item \thetodocounter:} #2}}
\newtheorem{theorem}{Theorem}
\newtheorem{theorem2}{Theorem}
\newtheorem{theorem3}{Theorem}
\newtheorem{theorem4}{Theorem}
\newtheorem{theorem5}{Theorem}
\newtheorem{acknowledgement}[theorem]{Acknowledgement}
\newtheorem{algorithm}[theorem]{Algorithm}
\newtheorem{axiom}[theorem]{Axiom}
\newtheorem{case}[theorem]{Case}
\newtheorem{claim}[theorem]{Claim}
\newtheorem{conclusion}[theorem]{Conclusion}
\newtheorem{condition}[theorem]{Condition}
\newtheorem{conjecture}[theorem4]{Conjecture}
\newtheorem{corollary}{Corollary}
\newtheorem{assumption}{Assumption}
\newtheorem{criterion}[theorem]{Criterion}
\newtheorem{definition}[theorem]{Definition}
\newtheorem{example}[theorem]{Example}
\newtheorem{exercise}[theorem]{Exercise}
\newtheorem{lemma}{Lemma}
\newtheorem{notation}[theorem]{Notation}
\newtheorem{problem}[theorem]{Problem}
\newtheorem{proposition}[theorem]{Proposition}
\newtheorem{remark}[theorem]{Remark}
\newtheorem{solution}[theorem]{Solution}
\newtheorem{summary}[theorem]{Summary}

%%%%%%%%%%%%%%%%%%%%%%%%%%%%%%%%%%%%%%%%%%%%%%%%%%%%%%%%%%%%%%%%%
%%%%%%%%%%%%%%%%%%%%%%%%%%%%%%%%%%%%%%%%%%%%%%%%%%%%%%%%%%%%%%%%%
%%%%%%%%%%%%%%%%%%%%%%%%%%%%%%%%%%%%%%%%%%%%%%%%%%%%%%%%%%%%%%%%%
%The flowchart stuff
\tikzstyle{att0} = [rectangle, rounded corners, 
minimum width=2cm, 
minimum height=1cm,
text centered, 
draw=black, 
fill=red!30]

\tikzstyle{att1} = [rectangle, rounded corners, 
minimum width=2cm, 
minimum height=1cm,
text centered, 
draw=black, 
fill=blue!30]

\tikzstyle{arrow} = [thick,->,>=stealth]

\addbibresource{refs.bib}
